\documentclass[10pt,a4paper]{amsart}
\usepackage[margin=19mm]{geometry}
\usepackage{amsmath,amssymb,mathtools}
\usepackage[scr=rsfs,cal=euler]{mathalfa}
\usepackage{xfrac}
\usepackage[unicode,hidelinks,bookmarks=false]{hyperref}
\newcommand{\ie}{i.e.\ }
\newcommand{\cf}{cf.\ }
\newcommand{\resp}{respectively\ }
\newcommand{\Subsection}{\subsection}
\providecommand{\nobreakdash}{\nobreak\hbox{-}}
\setlength{\emergencystretch}{2em}
\multlinegap0pt
\usepackage{amssymb}
\usepackage{mathtools}
\newtheorem{lemma}{Lemma}
\newcommand{\h}{h_2}
\title{\textbf{Finite-Degree Quantum LDPC Codes Reaching the Gilbert--Varshamov Bound}}
\author{Kenta Kasai}
\date{Source: arXiv:2603.24588v3 (CC BY 4.0)}
\begin{document}
\maketitle

\noindent\emph{Source and licence.} \textbf{Finite-Degree Quantum LDPC Codes Reaching the Gilbert--Varshamov Bound}. Version arXiv:2603.24588v3 (CC BY 4.0), distributed by arXiv under the Creative Commons Attribution 4.0 International licence, http://creativecommons.org/licenses/by/4.0/, as recorded in the arXiv licence metadata for this exact version and shown on the abstract page https://arxiv.org/abs/2603.24588v3. Full licence notice: \texttt{licenses/arxiv-2603.24588-CC-BY-4.0.txt}.\par\smallskip

\noindent\textit{Editorial scope.} Counted unit: the article's lemma lem:mn-window-bounds-jp (source0.tex lines 2734--2771). The article's complete original proof of it is delivered with it (source0.tex lines 2773--2946, one \texttt{proof} environment of 174 lines). No mathematics is written, repaired or re-derived: the statement and the proof are the article's own lines, in the article's own notation, and no retained line is substituted or re-flowed.  Retained uncounted as the source-local context the proof needs: lines 1464--1471.  Nothing was omitted from any retained span, so the package carries no omission record.  No retained line contains a citation, so the wrapper carries no bibliography and no bibliography entry.  No retained line contains a figure environment, an image-inclusion command or drawing code, so no image is delivered and the package declares no asset dependency.  Editorial formatting only, in addition: an AMS \texttt{amsart} preamble that restates the article's own declarations and macros used by the retained lines, namely the environments \texttt{lemma} on separate counters, together with the standard packages those lines need.  The retained environments print in this document as Lemma 1 in order of appearance, whereas the article prints the same items with the numbering of its own full document.  The complete native source of the article as distributed in the arXiv e-print for this exact version is bundled unchanged as \texttt{sources/197155/source0.tex}.
\begin{lemma}[Scalar zero-output gap at the split scale]\label{lem:scalar-zero-output-gap-jp}
Let \(k\ge10\), \(2\le L<k\), and \(\alpha:=L/k\). Set \(q_0:=2\mathrm{e}\,\mathrm{e}^{-k}\). Then
\[
\Psi_{L,k}(q):=
\alpha\h(q)-1+\log_2\!\bigl(1+(1-2q)^L\bigr)
\]
is strictly negative on \(q_0\le q\le1/2\).
\end{lemma}

\begin{lemma}[Elementary bounds for the finite degree window]\label{lem:mn-window-bounds-jp}
For all parameters above, the following statements hold.

\begin{enumerate}
\item Let
\[
\rho:=\max\!\left\{\frac\beta{j_Z},\frac\beta{j_\Delta},\omega_*\right\},
\quad c_X:=\frac{2\beta}{k}+\omega_*,
\quad C_X:=\frac{1+\rho}{(1-\rho)^2},
\]
and
\[
B_X:=\sqrt2\,\frac{\mathrm e(1+A)}{c_X}
\left(\frac{C_Xkc_X}{\mathrm e}\right)^{k/2}.
\]
Then \(B_X<7/20<1\).

\item Uniformly on
\[
\frac1{10L}\le q\le\frac12,\qquad 0\le\omega\le\omega_*,
\]
one has \(\psi(q,\omega)<0\).  In fact the upper bound obtained in the
proof below is at most \(-1/100\) in natural-log units.

\item If \(j_Z\) is odd, set
\[
\eta:=
\left(1-\frac{2\beta}{j_Z}\right)^{j_Z}
\left(1-\frac{2\beta}{j_\Delta}\right)^{j_\Delta}.
\]
Then
\[
\alpha_Z\h\!\left(\frac\beta{j_Z}\right)
+\alpha_\Delta\h\!\left(\frac\beta{j_\Delta}\right)
+\h(\omega_*)-1+\log_2\!\left(1-\frac A2\eta\right)<-\frac{23}{25}< -\frac9{10}.
\]
\end{enumerate}
\end{lemma}

\begin{proof}
We first state explicitly how the finitely many degree-dependent scalar
comparisons are made.  Write a positive rational argument as
\(x=2^py\), where \(p\in\mathbb Z\) and \(1\le y<2\), and use
\[
 \ln y=2\sum_{r=0}^{N-1}\frac{z^{2r+1}}{2r+1}+R_N,
 \quad z=\frac{y-1}{y+1},\qquad
 0\le R_N\le\frac{2z^{2N+1}}{(2N+1)(1-z^2)}            \tag{MN--1}
\]
with strict inequalities when \(y>1\).  Then
\(\ln x=p\ln2+\ln y\), using the same series at \(z=1/3\) for
\(\ln2\).  This gives rational upper and lower bounds.  The algebraic number
\(s=(A/2)^{1/k}\) is bracketed by adjacent rationals \(s^-<s<s^+\),
verified exactly by \((s^-)^k<A/2<(s^+)^k\); hence
\((1-s^+)/2<\omega_*<(1-s^-)/2\).  Taking \(N=72\) and denominator
\(10^{12}\) for the root brackets gives the following finite table.
The third column reports the case whose exact upper bound is largest, and
the last column gives the rational threshold actually used in the proof.
\[
\begin{array}{l|c|c|c}
\text{quantity}&\text{cases}&\text{case }(j_Z,L,k)&
 \text{proved upper bound}\\ \hline
B_X&56&(4,6,10)& <17/50\\
H(\omega_*)-\alpha_Z\ln2&56&(4,26,30)&<-19/1000\\
z_0-Ar_0^2/2-\alpha_Z\ln2+\ln(1+\mathrm e^{-c_0})
 &56&(4,26,30)&<-13/100\\
\text{odd-corner expression}&24&(5,7,12)&<-23/25
\end{array}                                                     \tag{MN--2}
\]
Here \(r_0,z_0,c_0\) are the quantities defined below.  For the first
row we use the rational inequalities
\(\sqrt2<70711/50000\) and \(\mathrm e>1359/500\); these follow by
squaring the first rational and by
\(\sum_{r=0}^8 1/r!>1359/500\), respectively.  In the third row,
\(\ln(1+\mathrm e^{-c_0})\le\mathrm e^{-c_0}<r\) is certified by the
rational test \(c_0+\ln r>0\).  Thus (MN--2) is a list of exact scalar
inequalities, rather than a floating-point maximization or an enclosure
of a continuous box.  The supplementary script
\texttt{verify\_analytic\_scalar\_bounds.py} implements only the fraction
operations in (MN--1), enumerates the 56 admissible pairs, and reproduces
the table.  The first and fourth rows prove assertions 1 and 3.

We prove the second assertion in detail.  First suppose
\(1/(10L)\le q\le1/(2L)\), and set \(x=Lq\),
\(z=(1-2q)^L\).  The function
\[
g(\omega):=\ln\!\left(1+z(1-2\omega)^k\right)
\]
is convex because
\[
g''(\omega)=
\frac{4kz(1-2\omega)^{k-2}
\bigl((k-1)-z(1-2\omega)^k\bigr)}
{\bigl(1+z(1-2\omega)^k\bigr)^2}\ge0.
\]
Consequently its chord on \([0,\omega_*]\) gives
\[
g(\omega)\le\ln(1+z)-c(q)\omega,
\quad
c(q):=\frac{\ln(1+z)-\ln(1+Az/2)}{\omega_*}.
\]
The entropy conjugacy formula
\[
\sup_{0\le w\le1}\{H(w)-cw\}=\ln(1+\mathrm e^{-c})
\]
therefore yields
\[
\psi(q,\omega)
\le F_0(q)+\ln(1+\mathrm e^{-c(q)}),
\quad
F_0(q):=AH(q)-\ln2+\ln(1+(1-2q)^L).
\]
As in Lemma~\ref{lem:scalar-zero-output-gap-jp}, for \(x\le1\),
\[
F_0(q)\le
f_{L,k}(x):=
x\left[\frac1k\ln\frac{\mathrm eL}{x}-1+\frac x2\right].
\]
The function \(f_{L,k}\) is convex on \([1/10,1/2]\).  Elementary
endpoint comparisons give
\[
\begin{array}{c|ccc}
x&1/10&1/4&1/2\\ \hline
\text{upper bound for }f_{L,k}(x)&-43/1000&-113/1000&-1/5
\end{array}
\]
Indeed, \(L\le k-4\), and
\([1+\ln((k-4)/x)]/k\) decreases for \(k\ge10\) and
\(x\le1/2\), so all three upper bounds reduce to \((L,k)=(6,10)\)
and follow from (MN--1).
Explicitly, the numerator of its derivative is
\[
 \frac{k}{k-4}-\left[1+\ln\frac{k-4}{x}\right]
 \le\frac53-[1+\ln12]<0.
\]

We next give uniform lower bounds for \(c(q)\), without a degree scan.
For every admissible triple,
\[
 \omega_*<\frac{57}{1000}.                                  \tag{MN--3}
\]
To see this, put \(s=1-2\omega_*\).  If \(k=10\), then
\(s^{10}=A/2\ge3/10>(443/500)^{10}\).  If \(k\ge12\), then
\(A/2\ge3/11>(443/500)^{12}\ge(443/500)^k\).  Here
\(A\ge6/11\) follows from
\(A\ge1/2+1/k\ge11/20>6/11\) for \(k\le20\), and from
\(A\ge1-10/k\ge6/11\) for \(k\ge22\).  Thus
\(s>443/500\) in both cases, proving (MN--3).

For fixed \(x=Lq\), the value
\(z=(1-2x/L)^L\) increases with \(L\): with \(y=2x/L\), the derivative
of its logarithm has sign
\(\ln(1-y)+y/(1-y)>0\).  Moreover,
\(\ln(1+z)-\ln(1+Az/2)\) increases with \(z\) and decreases with
\(A\); its derivative with respect to \(z\) is
\((1-A/2)/[(1+z)(1+Az/2)]>0\).  Since \(L\ge6\) and
\(A\le13/15\), (MN--1) gives
\[
c(1/(4L))>
\frac{\ln(1+(11/12)^6)-\ln(1+(13/30)(11/12)^6)}{57/1000}>4,
\]
and
\[
c(1/(2L))>
\frac{\ln(1+(5/6)^6)-\ln(1+(13/30)(5/6)^6)}{57/1000}>\frac52.
\]
The two logarithmic numerators are respectively larger than
\(237/1000\) and \(1533/10000\).  Because \(z(q)\) decreases and the numerator
of \(c(q)\) increases with \(z\), the function \(c(q)\) decreases with \(q\);
these estimates hold throughout the corresponding subintervals.
Thus on \(1/10\le x\le1/4\),
\[
\psi<-\frac{43}{1000}+\frac1{50}=-\frac{23}{1000},
\]
while on \(1/4\le x\le1/2\),
\[
\psi<-\frac{113}{1000}+\frac1{12}< -\frac{29}{1000}.
\]
Here \(\mathrm e^{-4}<1/50\) and \(\mathrm e^{-5/2}<1/12\), again
by (MN--1).  In particular both bounds are below \(-1/100\).

It remains to treat \(q\ge1/(2L)\).  Put
\(r=1-2q\), \(r_0=1-1/L\), and \(z_0=r_0^L\).  Pinsker's inequality gives
\(\ln2-H(q)\ge r^2/2\), and \(\ln(1+y)\le y\), so
\[
\psi(q,\omega)
\le H(\omega)-\alpha_Z\ln2
+(1-2\omega)^kr^L-\frac A2r^2.
\]
As a function of \(r^2\), the last two terms are convex because \(L\ge6\).
Their maximum on \([0,r_0^2]\) is therefore at an endpoint.  At \(r=0\)
the resulting bound is
\[
H(\omega_*)-\alpha_Z\ln2<-\frac{19}{1000}<-\frac1{100}
\]
by the second row of (MN--2).
At \(r=r_0\), convexity of \((1-2\omega)^k\) on
\([0,\omega_*]\) gives
\[
(1-2\omega)^k
\le1-\left(1-\frac A2\right)\frac\omega{\omega_*}.
\]
Writing
\[
c_0:=\frac{z_0(1-A/2)}{\omega_*},
\]
and applying entropy conjugacy once more bounds this endpoint by
\[
z_0-\frac A2r_0^2-\alpha_Z\ln2+\ln(1+\mathrm e^{-c_0})
<-\frac{13}{100}
\]
by the third row of (MN--2).
This proves the second assertion.
\end{proof}

\end{document}
